\section{4. Method}

\subsection{4.1 Overview}

Galerkin Attention has two execution paths:

\begin{itemize}
\item \textbf{Encoder path:} Petrov-Galerkin projection onto DCT basis, then attention
\end{itemize}
  at rank $r$.
\begin{itemize}
\item \textbf{Decoder path:} Standard causal softmax attention (fallback).
\end{itemize}

\textbf{Figure 4.1} shows the pipeline.

\begin{lstlisting}
    +-------------------------------------------+
    |  Input X $\to$ Q/K/V projections (frozen+LoRA)|
    +---------------------+---------------------+
                          |
                           
    +-------------------------------------------+
    |  Apply RoPE, GQA repeat_kv                |
    +---------------------+---------------------+
                          |
                           
    +-------------------------------------------+
    |  Is decoder?                              |
    +----+---------------------+----------------+
         | Yes                 | No
                                
    +--------------+    +----------------------+
    |  Standard    |    |  1. Build DCT basis  |
    |  causal      |    |  2. Project Q/K/V    |
    |  softmax     |    |  3. Attention in r   |
    |  attention   |    |  4. Reconstruct + $\sqrt{}$  |
    +------+-------+    +----------+-----------+
           |                       |
           +-----------+-----------+
                        
              +----------------+
              |  out_proj      |
              +--------+-------+
                       |
                        
                      Y
\end{lstlisting}

\subsection{4.2 Phase 1: Projections}

Frozen base + LoRA:

\[
\text{Proj}_P(X) = X W_P^\top + \frac{\alpha}{r_L} (X A_P^\top) B_P^\top + b_P
\tag{4.1}
\]

$W_P$ copied exactly from pretrained.

\subsection{4.3 Phase 2: Basis Construction}

\subsubsection{DCT-II Basis}

For sequence length $N$ and rank $r$:

\[
\Phi_{n, k} = \sqrt{\frac{2 - \delta_{k0}}{N}} \cos\!\left(\frac{\pi k (2n + 1)}{2N}\right)
\tag{4.2}
\]

with indices $n = 0, \dots, N-1$ and $k = 0, \dots, r-1$.

\subsubsection{Orthonormalization}

DCT is theoretically orthonormal, but floating-point rounding can break 
this. We add small identity regularization and apply QR:

\[
\Phi \leftarrow \text{QR}(\Phi + \epsilon I_{N \times r})
\tag{4.3}
\]

with $\epsilon = 10^{-4}$.

\subsection{4.4 Phase 3: Encoder Path}

\subsubsection{Projection to Basis}

\[
Q_\Phi = \Phi^\top Q \in \mathbb{R}^{B \times H \times r \times d_h}
\tag{4.4}
\]

\[
K_\Phi = \Phi^\top K_{\text{masked}}, \quad V_\Phi = \Phi^\top V_{\text{masked}}
\tag{4.5}
\]

where $K_{\text{masked}} = K \cdot m$ and $V_{\text{masked}} = V \cdot m$ 
with padding mask $m$.

\subsubsection{Attention in Basis Space}

\[
\text{Scores} = \frac{Q_\Phi K_\Phi^\top}{\sqrt{d_h}} \in \mathbb{R}^{B \times H \times r \times r}
\tag{4.6}
\]

\[
\text{Probs} = \text{softmax}(\text{Scores})
\tag{4.7}
\]

\[
O_\Phi = \text{Probs} \cdot V_\Phi
\tag{4.8}
\]

\subsubsection{Reconstruction to Sequence Grid}

\[
Y = \Phi \cdot O_\Phi \cdot \sqrt{N / r}
\tag{4.9}
\]

The scaling factor $\sqrt{N / r}$ compensates for the energy loss due to 
truncation.

\subsubsection{Padding Handling}

Masked tokens in $K$ and $V$ are zeroed before projection. After 
reconstruction, the output at padded positions is zeroed:

\[
Y \leftarrow Y \odot m
\tag{4.10}
\]

\subsection{4.5 Phase 4: Decoder Path}

For causal models, projection breaks causality because $K_\Phi$ depends on 
\textbf{all} keys. We fall back to standard softmax attention:

\[
\text{Scores} = \frac{Q K^\top}{\sqrt{d_h}}, \quad
\text{Scores} \leftarrow \text{Scores} + M_{\text{causal}}
\tag{4.11}
\]

\[
\text{Probs} = \text{softmax}(\text{Scores}), \quad
Y = \text{Probs} \cdot V
\tag{4.12}
\]

with numerical stabilization via $-\max$ subtraction.

\subsection{4.6 The Full Algorithm}

\textbf{Algorithm 1:} Galerkin Attention forward
\begin{lstlisting}
Input:  X, attention_mask, position_embeddings
        is_decoder, rank r
Output: Y, probs

1:  Q $\leftarrow$ LoRALinear_Q(X)
2:  K $\leftarrow$ LoRALinear_K(X)
3:  V $\leftarrow$ LoRALinear_V(X)
4:
5:  if RoPE: (Q, K) $\leftarrow$ apply_rope(Q, K, cos, sin)
6:  if GQA: K, V $\leftarrow$ repeat_kv(K, V)
7:
8:  if is_decoder:                                DECODER FALLBACK
9:      S $\leftarrow$ Q K^T / $\sqrt{}$d_h
10:     S $\leftarrow$ S + causal_mask
11:     S $\leftarrow$ S + padding_mask
12:     P $\leftarrow$ softmax(S)
13:     Y_ctx $\leftarrow$ P @ V
14:
15: else:                                          ENCODER GALERKIN
16:     mask $\leftarrow$ valid_token_mask(attention_mask)
17:     K_masked $\leftarrow$ K   mask
18:     V_masked $\leftarrow$ V   mask
19:
20:       $\leftarrow$ DCT_basis(N, r)                          Orthonormalize via QR
21:     Q_  $\leftarrow$  ^T Q
22:     K_  $\leftarrow$  ^T K_masked
23:     V_  $\leftarrow$  ^T V_masked
24:
25:     S $\leftarrow$ Q_  K_ ^T / $\sqrt{}$d_h                         (B, H, r, r)
26:     P $\leftarrow$ softmax(S)
27:     O_  $\leftarrow$ P @ V_ 
28:
29:     Y_ctx $\leftarrow$   @ O_                                Reconstruct
30:     Y_ctx $\leftarrow$ Y_ctx   $\sqrt{}$(N/r)                        Energy correction
31:     Y_ctx $\leftarrow$ Y_ctx   mask                          Zero padded positions
32:
33: Y_ctx $\leftarrow$ reshape(B, N, H   d_h)
34: if nested: Y $\leftarrow$ Y_ctx
35: else: Y $\leftarrow$ LoRALinear_O(Y_ctx)
36:
37: return Y
\end{lstlisting}

\subsection{4.7 Complexity Analysis}

\begin{table}[H]
\centering
\begin{tabular}{lll}
\toprule
\textbf{Path} & \textbf{Time} & \textbf{Memory} \\
\midrule
Encoder (Galerkin) & $\mathcal{O}(N r d)$ + $\mathcal{O}(N r)$ for basis & $\mathcal{O}(N r)$ \\
Decoder (fallback) & $\mathcal{O}(N^2 d)$ & $\mathcal{O}(N^2)$ \\
\bottomrule
\end{tabular}
\end{table}

For $N = 128$ and $r = 64$:
\begin{itemize}
\item Encoder: $\sim 2 \times 10^6$ FLOPs per head
\end{itemize}
\begin{itemize}
\item Decoder: $\sim 4 \times 10^6$ FLOPs per head
\end{itemize}
\begin{itemize}
\item \textbf{Encoder speedup: ~2$\times$}
\end{itemize}

For $N = 512$ and $r = 64$:
\begin{itemize}
\item Encoder: $\sim 8 \times 10^6$
\end{itemize}
\begin{itemize}
\item Decoder: $\sim 6.4 \times 10^7$
\end{itemize}
\begin{itemize}
\item \textbf{Encoder speedup: ~8$\times$}
\end{itemize}

\subsection{4.8 Parameter Analysis}

\begin{table}[H]
\centering
\begin{tabular}{lll}
\toprule
\textbf{Component} & \textbf{Baseline} & \textbf{Galerkin} \\
\midrule
Base weights & $4 d^2$ & $4 d^2$ (frozen) \\
LoRA ($r_L$) & 0 & $8 r_L d$ \\
**Total per layer** & $4 d^2$ & $4 d^2 + 8 r_L d$ \\
\bottomrule
\end{tabular}
\end{table}

For DistilBERT ($d = 768$, $r_L = 16$): \textbf{+589,824 params} across 6 layers.

\textbf{Galerkin Attention is not a compression method.}

\subsection{4.9 Theoretical Notes}

\subsubsection{4.9.1 Galerkin Approximation Error}

\textbf{Claim.} The Galerkin approximation error for the attention operator is 
bounded by the DCT tail energy:

\[
\|\mathcal{K} Q - \hat{\mathcal{K}} \hat{Q}\|_2 \leq C \cdot \sqrt{\sum_{k > r} \sigma_k^2}
\tag{4.13}
\]

where $\sigma_k$ are the singular values of the kernel operator.

\textbf{Proof sketch.} Standard Galerkin theory (Strang, 1973, Theorem 5.2). 
The error is dominated by the tail of the eigen-decomposition of $\mathcal{K}$. 
DCT captures the low-frequency modes optimally. $\square$

\subsubsection{4.9.2 Energy Preservation Scaling}

\textbf{Claim.} The scaling $\sqrt{N/r}$ preserves the expected output energy 
under the assumption that $O_\Phi$ coefficients are i.i.d.

\textbf{Proof.} For $O_\Phi \in \mathbb{R}^{r \times d_h}$ with i.i.d. entries 
variance $\sigma^2$, the reconstruction $Y = \Phi O_\Phi$ has entries:

\[
Y_{n, j} = \sum_k \Phi_{n,k} (O_\Phi)_{k, j}
\tag{4.14}
\]

with variance:

\[
\text{Var}(Y_{n,j}) = \sum_k \Phi_{n,k}^2 \sigma^2 = \sigma^2 \cdot \|\Phi_{n,:}\|_2^2 = \sigma^2
\tag{4.15}
\]

since $\Phi$ has orthonormal columns. The output energy per sequence 
position is thus $\sigma^2$, matching the expected basis energy $\sigma^2 \cdot r$. 
The $\sqrt{N/r}$ factor compensates for the fact that $r$ basis modes must 
represent $N$ sequence positions. $\square$

\subsubsection{4.9.3 Causality Preservation for Encoders}

Encoders have no causal constraint. The Galerkin projection preserves 
bidirectional context. $\square$

\subsection{4.10 Hyperparameter Guidance}
